%% ch05.tex — Chapter 5: Resource Management, File Formats, and I/O Optimisation
%% Part of "Data Engineering" textbook

\chapter{Resource Management, File Formats, and I/O Optimisation}
\label{ch:yarn-formats}

\epigraph{\textit{``Storing data in HDFS is cheap. Storing it in the wrong format
is expensive --- you pay the cost every single time you query it.''}}{Julien Le Dem, Apache Parquet co-creator, 2014}

\vspace{4pt}

\begin{chapterlearning}
After completing this chapter, you will be able to:
\begin{enumerate}
  \item Identify the specific bottlenecks of the Hadoop 1.x JobTracker
        architecture and explain why they imposed an approximate
        4{,}000-node scalability ceiling.
  \item Describe YARN's three-component architecture---Resource\-Manager,
        NodeManager, and ApplicationMaster---and trace the complete
        job execution flow from client submission to container completion.
  \item Explain what a YARN Container is, why it replaced the fixed
        MapReduce \emph{slot} model, and how it enables framework-agnostic
        resource sharing.
  \item Compare YARN's three pluggable schedulers---FIFO, Capacity, and
        Fair---and select the appropriate scheduler for a given
        multi-tenant production scenario.
  \item Describe YARN's three-tier fault tolerance model (NodeManager,
        ApplicationMaster, and ResourceManager failures) and identify
        which tier each failure type is handled at.
  \item Explain the fundamental difference between row-oriented and
        columnar storage, derive the I/O reduction formula for a
        $k$-column query on a $c$-column table, and apply it to a
        concrete performance calculation.
  \item Describe the internal structure of a Parquet file---row groups,
        column chunks, and pages---and explain how the footer enables
        predicate pushdown to skip row groups without reading them.
  \item Compare Avro, Parquet, and ORC on orientation, schema model,
        splittability, compression, and primary use case; and select
        the appropriate format for a given data engineering workload.
  \item Explain the splittability property for compression codecs and
        explain why Gzip-compressed CSV is an anti-pattern in HDFS.
  \item Describe Dominant Resource Fairness (DRF) and explain why
        single-resource fairness metrics fail for heterogeneous workloads.
\end{enumerate}
\end{chapterlearning}

\bigskip

Chapters~\ref{ch:hdfs} and~\ref{ch:mapreduce} established the two
pillars of the original Hadoop stack: HDFS stores data, and MapReduce
processes it. By 2012, however, two engineering problems were threatening
the viability of that stack at production scale.

The first problem was \emph{scheduling}. A Hadoop 1.x cluster could only
run MapReduce jobs; when Apache Spark (2010) and Apache Tez (2013)
emerged as significantly faster alternatives, they had no path to cluster
resources because the MapReduce JobTracker was hardcoded to manage only
MapReduce tasks. Organisations that wanted to use Spark had to maintain
a completely separate cluster, with all the associated duplication of
hardware, operations effort, and data movement overhead.

The second problem was \emph{storage format}. As analytical workloads
evolved from full-table scans to targeted column selections---driven by
the growth of SQL analytics (Hive) and machine learning feature
engineering---the row-oriented file formats that Hadoop had inherited
(CSV, JSON, sequence files) became serious performance bottlenecks.
Selecting 3 columns from a 100-column CSV table on HDFS required reading
and discarding 97\% of the data on every query.

YARN \parencite{YARN2013} and columnar file formats (Parquet, ORC)
were the industry's responses to these two problems. This chapter covers
both in depth, then examines the compression and I/O optimisation
techniques that complete the performance picture.

%%====================================================================
\section{The Hadoop 1.x Scheduling Bottleneck}
\label{sec:hadoop1-bottleneck}
%%====================================================================

In Hadoop 1.x, cluster resource management and MapReduce job execution
were fused into a single process: the \term{JobTracker}. The JobTracker
performed four distinct functions simultaneously:
\begin{enumerate}
  \item Accepted job submissions from clients.
  \item Allocated fixed \term{slots}---Map slots and Reduce slots---to
        tasks running on TaskTrackers.
  \item Monitored every running task and detected failures.
  \item Maintained the job history for completed jobs.
\end{enumerate}

This fusion created three interlocking problems that collectively
prevented Hadoop 1.x from scaling beyond approximately 4{,}000 nodes.

\textbf{Framework lock-in.} The slot allocation model was hardcoded for
MapReduce. A Map slot could run exactly one Map task; a Reduce slot could
run exactly one Reduce task. There was no abstraction for the resource
requirements of a non-MapReduce task---no way for a Spark executor to
request a container of 8\GB\ RAM and 4 CPU cores, for example. Any
framework that needed to run on the cluster had to pretend to be MapReduce.

\textbf{Rigid slot granularity.} Map slots and Reduce slots were
configured with fixed memory and CPU limits cluster-wide. A job with
small Map tasks and large Reduce tasks could not allocate more memory
to Reducers at the expense of Mappers; both were constrained by the
same slot size. A node with 10 Map slots and 0 available Reduce slots
could not repurpose its unused Map slots to serve waiting Reduce tasks,
even if it had spare RAM and CPU capacity.

\textbf{Single-point scalability.} The JobTracker's memory footprint
grew linearly with the number of running tasks: each task required a
metadata entry in the JobTracker's heap. At 4{,}000 nodes with 10 tasks
per node, the JobTracker managed 40{,}000 task records simultaneously.
Beyond this scale, JVM garbage collection pressure made the JobTracker
unresponsive \parencite{YARN2013}.

%%====================================================================
\section{YARN: Yet Another Resource Negotiator}
\label{sec:yarn}
%%====================================================================

Apache Hadoop YARN, published by Vavilapalli et al.\ in 2013, resolved
all three bottlenecks through a single architectural decision: separate
cluster resource management from application-specific execution logic
\parencite{YARN2013}. YARN replaces the monolithic JobTracker with three
distinct components, each with a clearly bounded responsibility.

%--------------------------------------------------------------------
\subsection{The ResourceManager}
\label{subsec:rm}
%--------------------------------------------------------------------

The \term{ResourceManager} (RM) is the cluster-level resource scheduler.
It runs on a dedicated master node and contains two sub-components.

The \term{Scheduler} is responsible for allocating \term{containers}---
the YARN resource abstraction---to running applications. The Scheduler
is \emph{pluggable}: operators can choose from three built-in policies
(described in Section~\ref{sec:schedulers}) or implement a custom policy.
Critically, the Scheduler has no knowledge of the application logic; it
only knows the container requirements (\texttt{\{vcores: N, memory:
M~MB\}}) and the cluster's current capacity. This separation means the
same Scheduler can manage MapReduce tasks, Spark executors, Samza stream
processors, and any other future framework without modification.

The \term{ApplicationsManager} accepts job submissions from clients and
coordinates the launch of each application's first container, which will
host the ApplicationMaster.

%--------------------------------------------------------------------
\subsection{The NodeManager}
\label{subsec:nm}
%--------------------------------------------------------------------

A \term{NodeManager} (NM) runs on every worker node in the cluster.
Its responsibilities are narrowly scoped to the local machine:
\begin{itemize}
  \item Launch containers requested by ApplicationMasters, using Linux
        cgroups for CPU isolation and enforced memory limits.
  \item Monitor the resource usage (CPU, RAM, disk I/O, network) of
        every container on the local node and report to the ResourceManager
        via periodic heartbeats.
  \item Kill any container that exceeds its allocated memory limit to
        protect other containers on the same node.
  \item Aggregate container logs and upload them to HDFS for post-job
        debugging.
\end{itemize}

The NodeManager has no global view of the cluster; it knows only about
its own node. This is a deliberate design choice: the scalability of the
YARN control plane depends on the RM and each NM having well-bounded
state, with the RM collecting aggregate information rather than
processing per-task details.

%--------------------------------------------------------------------
\subsection{The ApplicationMaster}
\label{subsec:am}
%--------------------------------------------------------------------

The \term{ApplicationMaster} (AM) is the component that makes YARN
framework-agnostic. One ApplicationMaster runs per submitted job, inside
its own container allocated by the ResourceManager. The AM contains all
application-specific logic that the old JobTracker had embedded:

\begin{itemize}
  \item \textbf{Resource negotiation:} the AM communicates with the RM to
        request containers for its tasks. A Spark AM requests containers
        with specific CPU and memory configurations for Spark Executors;
        a MapReduce AM requests separate Map-phase and Reduce-phase
        containers.
  \item \textbf{Task placement:} the AM instructs NodeManagers to launch
        specific tasks inside allocated containers. The AM can apply its
        own data-locality preferences when choosing which nodes to request
        containers from.
  \item \textbf{Progress monitoring:} the AM tracks the progress of all
        running tasks and re-requests containers for failed tasks.
  \item \textbf{Framework-specific fault tolerance:} the MapReduce AM
        implements speculative execution; the Spark AM manages executor
        loss and task re-submission.
\end{itemize}

The AM model means YARN is, in the authors' words, ``a cluster operating
system'' \parencite{YARN2013}: the RM and NMs are the kernel that
manages hardware resources; the AM is a user-space process that
implements the policy of a specific framework. Just as Linux supports
arbitrary user-space processes, YARN supports arbitrary AMs.

%--------------------------------------------------------------------
\subsection{The Container: YARN's Resource Unit}
\label{subsec:container}
%--------------------------------------------------------------------

\begin{defbox}{YARN Container}
A \textbf{YARN Container} is an allocation of a specific quantity of
cluster resources on a specific node:
\[
\text{Container} = \bigl\{\;\text{node} : m,\;\; \text{vCores} : n,\;\;
                           \text{memory} : M\;\text{MB}\;\bigr\}
\]
The Container replaces Hadoop 1.x's fixed Map/Reduce slot with a
flexible, application-specified resource envelope. The NodeManager
on node $m$ enforces the allocated limits; any container exceeding
its memory allocation is immediately killed by the NM.
\end{defbox}

The Container abstraction is the key enabler of framework diversity.
Because a Container is specified in terms of generic resources (vCores
and RAM)---not in terms of task types---the same Resource\-Manager can
grant containers to a MapReduce AM requesting small (2~vCPU,\,4~GB)
Map-task containers, a Spark AM requesting large (8~vCPU,\,32~GB)
executor containers, and a Samza AM requesting tiny (1~vCPU,\,2~GB)
stream-task containers, all simultaneously on the same cluster.

%--------------------------------------------------------------------
\subsection{YARN Job Execution Flow}
\label{subsec:yarn-flow}
%--------------------------------------------------------------------

The YARN job execution follows six steps from client submission to
completion:

\begin{enumerate}
  \item \textbf{Submit.} The client submits the job to the
        ResourceManager's ApplicationsManager, providing the application
        JAR (or binary), configuration, and resource requirements.

  \item \textbf{Launch AM container.} The RM allocates one container on
        an available NodeManager for the ApplicationMaster. The NM
        launches the AM process inside that container.

  \item \textbf{AM registers and requests resources.} The AM registers
        itself with the RM and submits resource requests: a list of
        containers with their required vCores, RAM, and preferred nodes
        (for data locality).

  \item \textbf{RM grants containers.} The Scheduler allocates containers
        based on the scheduling policy and current cluster availability.
        Granted containers are communicated to the AM.

  \item \textbf{AM launches tasks.} For each granted container, the AM
        contacts the appropriate NodeManager via RPC and instructs it to
        launch a specific task process inside the container.

  \item \textbf{Task execution, monitoring, and completion.} Tasks run
        inside containers, report progress to the AM. On task completion,
        the AM releases the container back to the RM. When all tasks
        complete, the AM unregisters from the RM, and the RM releases
        the AM's own container.
\end{enumerate}

\begin{figure}[ht]
\centering
\begin{tikzpicture}[
    actor/.style={rectangle, draw=DEBlue!60, fill=DELightBlue!40, rounded corners=4pt,
                  minimum width=2.6cm, minimum height=0.75cm, align=center,
                  font=\small\bfseries\sffamily},
    nm/.style={rectangle, draw=DEGreen!70!black, fill=DELightGreen, rounded corners=4pt,
               minimum width=2.2cm, minimum height=0.7cm, align=center,
               font=\small\sffamily},
    ambox/.style={rectangle, draw=DEAccent!80, fill=DELightAccent, rounded corners=4pt,
                  minimum width=2.2cm, minimum height=0.7cm, align=center,
                  font=\small\sffamily},
    arr/.style={-{Stealth[length=4pt,width=4pt]}, thick},
    lbl/.style={font=\scriptsize\sffamily, align=center}
  ]

  % Actors
  \node[actor] (cl)  at (0,   6.5) {Client};
  \node[actor] (rm)  at (5.0, 6.5) {Resource\-Manager};
  \node[ambox] (am)  at (5.0, 4.0) {Application\-Master};
  \node[nm]    (nm1) at (1.5, 1.5) {NodeManager 1};
  \node[nm]    (nm2) at (5.0, 1.5) {NodeManager 2};
  \node[nm]    (nm3) at (8.5, 1.5) {NodeManager 3};

  % Step 1: Client -> RM
  \draw[arr, DEAccent]
    (cl.east) -- (rm.west)
    node[midway, above, lbl, text=DEAccent] {\tiny \textbf{(1)} submit job};

  % Step 2: RM -> NM1 (launch AM container)
  \draw[arr, DEBlue]
    (rm.south) to[out=-100, in=90] (nm1.north)
    node[pos=0.4, left, lbl, text=DEBlue] {\tiny \textbf{(2)} launch\\AM container};

  % Step 3: AM -> RM (register + request)
  \draw[arr, DEAccent]
    (am.north) -- (rm.south)
    node[midway, right, lbl, text=DEAccent]
         {\tiny \textbf{(3)} register\\\tiny + request containers};

  % Step 4: RM -> AM (grant containers) - dashed
  \draw[arr, DEBlue, dashed]
    (rm.south) to[bend right=8] (am.north)
    node[midway, left, lbl, text=DEBlue]
         {\tiny \textbf{(4)} grant\\container list};

  % Step 5: AM -> NMs (launch tasks)
  \draw[arr, DEGreen!70!black]
    (am.south) to[out=-130, in=90] (nm1.north)
    node[pos=0.35, left, lbl, text=DEGreen!70!black] {\tiny \textbf{(5)} launch task};
  \draw[arr, DEGreen!70!black]
    (am.south) -- (nm2.north)
    node[midway, right, lbl, text=DEGreen!70!black] {\tiny \textbf{(5)} launch task};
  \draw[arr, DEGreen!70!black]
    (am.south) to[out=-50, in=90] (nm3.north)
    node[pos=0.35, right, lbl, text=DEGreen!70!black] {\tiny \textbf{(5)} launch task};

  % Step 6: NMs -> AM (task status)
  \draw[arr, DEGray!60, dashed]
    (nm2.north) to[bend right=15] (am.south)
    node[midway, left, lbl, text=DEGray!60] {\tiny \textbf{(6)} status/progress};

\end{tikzpicture}
\caption{YARN job execution flow. The ResourceManager (RM) manages global
         resources and launches the ApplicationMaster (AM) in its own container.
         The AM then negotiates additional containers with the RM and instructs
         NodeManagers to launch application tasks, monitoring their progress
         independently of the RM.}
\label{fig:yarn-flow}
\end{figure}

%%====================================================================
\section{YARN Schedulers: Sharing the Cluster}
\label{sec:schedulers}
%%====================================================================

With multiple teams submitting jobs simultaneously to a shared cluster,
the Scheduler's policy determines which jobs get resources first and
how much capacity each team is guaranteed. YARN provides three built-in
scheduler implementations.

%--------------------------------------------------------------------
\subsection{The FIFO Scheduler}
\label{subsec:fifo-sched}
%--------------------------------------------------------------------

The \term{FIFO (First In, First Out) Scheduler} maintains a single job
queue and assigns resources in submission order. The first job submitted
receives all available cluster capacity until it completes; subsequent
jobs wait. FIFO is simple to understand and configure, but it is
unsuitable for multi-tenant clusters. A single long-running batch job
(e.g., a multi-hour ETL pipeline) can starve time-sensitive interactive
queries for hours. FIFO is appropriate only for single-user development
clusters or automated pipeline clusters that run one job at a time.

%--------------------------------------------------------------------
\subsection{The Capacity Scheduler}
\label{subsec:capacity-sched}
%--------------------------------------------------------------------

The \term{Capacity Scheduler}, the default in Hadoop distributions from
Hortonworks and Cloudera, divides the cluster's total capacity into
named \term{queues}. Each queue has a guaranteed minimum capacity and a
maximum capacity:

\begin{itemize}
  \item The \textbf{minimum capacity} ensures that jobs in this queue can
        always obtain at least $X$\% of cluster resources, even when other
        queues are busy.
  \item The \textbf{maximum capacity} prevents a single queue from
        consuming the entire cluster, even when the cluster is otherwise idle.
  \item \textbf{Elastic borrowing}: when a queue has unused capacity, other
        queues can borrow it. The capacity is reclaimed---with a configurable
        grace period for preemption---when the owning queue submits new jobs.
\end{itemize}

A typical multi-team Hadoop cluster might configure:
\begin{center}
\begin{tabular}{lcc}
\toprule
\textbf{Queue}     & \textbf{Min capacity} & \textbf{Max capacity} \\
\midrule
\texttt{production} & 60\%                 & 80\% \\
\texttt{analytics}  & 30\%                 & 60\% \\
\texttt{adhoc}      & 10\%                 & 40\% \\
\bottomrule
\end{tabular}
\end{center}

The production queue always has at least 60\% of cluster capacity
available for revenue-generating pipelines, even if the analytics and
adhoc queues are idle.

%--------------------------------------------------------------------
\subsection{The Fair Scheduler}
\label{subsec:fair-sched}
%--------------------------------------------------------------------

The \term{Fair Scheduler}, the default in Apache Hadoop distributions and
historically used by Facebook and Twitter, allocates resources so that
all running jobs receive an equal share of available capacity over time.
When only one job is running, it receives 100\% of the cluster. When a
second job is submitted, it begins receiving resources immediately; over
a short window, both jobs converge on 50\% each. When the second job
completes, the first job again receives 100\%.

The key property of the Fair Scheduler is its treatment of newly
submitted short jobs. Under the Capacity Scheduler, a short interactive
query submitted to a busy production queue must wait behind already-running
batch jobs. Under the Fair Scheduler, the new query preempts some
resources from running jobs immediately, completing quickly before
returning those resources to the batch jobs. This makes the Fair Scheduler
particularly effective for clusters that serve a mix of batch pipelines
and interactive analytical queries.

%--------------------------------------------------------------------
\subsection{Dominant Resource Fairness}
\label{subsec:drf}
%--------------------------------------------------------------------

Both the Capacity and Fair Schedulers must decide how to measure
``fairness'' in a cluster with multiple resource dimensions (vCores and
RAM). A job requiring many CPUs but little memory and a job requiring
few CPUs but much memory have different resource profiles; allocating
equal CPU to both may leave one job RAM-starved while the other has
excess RAM.

\term{Dominant Resource Fairness} (DRF), introduced by Ghodsi et al.\
at NSDI 2011 \parencite{DRF2011}, resolves this by defining each job's
fairness share in terms of its \emph{dominant resource}---the resource
type for which the job's relative allocation is highest. If a job consumes
$20\%$ of the cluster's CPU but only $5\%$ of its RAM, its dominant
resource is CPU and its dominant share is $20\%$. The DRF scheduler
equalises dominant shares across all running jobs, ensuring that no job
dominates the cluster in the resource dimension that matters most to it.
YARN's Fair Scheduler implements DRF as its fairness metric for
multi-resource scheduling.

%--------------------------------------------------------------------
\subsection{YARN Fault Tolerance}
\label{subsec:yarn-fault}
%--------------------------------------------------------------------

YARN's fault tolerance is layered across three tiers, matching the
three components of its architecture.

\textbf{NodeManager failure.} If a NodeManager stops sending heartbeats
to the ResourceManager (default timeout: 10 minutes), the RM marks it
as dead. All containers on that NodeManager are immediately marked as
failed, and the RM notifies the relevant ApplicationMasters. Each AM
decides how to respond according to its framework's fault tolerance
policy: a MapReduce AM re-requests replacement containers for failed Map
or Reduce tasks; a Spark AM marks the executor as lost and reschedules
its tasks on surviving executors.

\textbf{ApplicationMaster failure.} If the AM process crashes, the RM
detects the failure via a missed heartbeat and restarts the AM in a
new container (up to a configurable maximum number of attempts, default
2). The restarted AM can recover job state from HDFS checkpoints if the
framework supports it. MapReduce AMs checkpoint completed task outputs
to HDFS, allowing a restarted AM to resume from the last completed
stage. Early versions of Spark did not checkpoint AM state, so AM
failure required restarting the entire job.

\textbf{ResourceManager failure.} In Hadoop 1.x, RM failure was
catastrophic: the entire cluster became inaccessible. YARN 2.4
introduced \term{ResourceManager High Availability}, mirroring the HDFS
HA design: an Active RM and a Standby RM, with ZooKeeper providing
leader election and automatic failover in under 60 seconds. The Active
RM's state (running applications, container allocations) is periodically
checkpointed to ZooKeeper so the Standby can reconstruct it on promotion.

\begin{researchspotlight}{Vavilapalli et al.\ (2013) --- Apache Hadoop YARN: Yet Another Resource Negotiator}
\textbf{Citation:} Vavilapalli, V.~K., et~al.\ (2013). Apache Hadoop
YARN: Yet Another Resource Negotiator. \textit{Proceedings of the
4th ACM Symposium on Cloud Computing (SoCC 2013)},
pp.\ 5:1--5:16. \parencite{YARN2013}

\smallskip\noindent
\textbf{Key insight:} The paper's central thesis is that the fusion of
resource management and MapReduce execution logic in Hadoop 1.x's
JobTracker was the primary scalability and extensibility bottleneck.
Separating these concerns---through the RM/NM/AM decomposition---makes
the cluster a general-purpose computing platform rather than a
MapReduce appliance.

\smallskip\noindent
\textbf{Technical contributions:}
\begin{itemize}
  \item Defined the Container abstraction (\texttt{\{node, vCores, memory\}})
        as the universal resource unit, replacing the rigid Map/Reduce
        slot model and enabling flexible, framework-specific resource
        allocation.
  \item Introduced the ApplicationMaster pattern: a per-application
        process that encapsulates framework-specific scheduling, fault
        tolerance, and monitoring logic, making YARN framework-agnostic
        by design.
  \item Demonstrated scalability to 6{,}000 nodes with 100{,}000
        simultaneous tasks, exceeding Hadoop 1.x's 4{,}000-node limit
        by 50\%, with ResourceManager scheduling latency below
        1 second at 1{,}000 heartbeats/second.
\end{itemize}

\smallskip\noindent
\textbf{Impact today:} YARN is cited over 3{,}000 times and became
the default cluster OS for the entire Hadoop ecosystem. Its three-tier
architecture directly influenced Kubernetes's Pod/Node/Master model;
the Container abstraction is the direct ancestor of the Kubernetes Pod.
Apache Spark, Apache Tez, Apache Samza, Apache Flink, and TensorFlow
on Hadoop all ran as YARN applications. The ApplicationMaster pattern
established a design precedent that every subsequent cluster computing
framework has replicated.
\end{researchspotlight}

%%====================================================================
\section{Row-Oriented vs.\ Columnar Storage}
\label{sec:row-columnar}
%%====================================================================

The second major performance dimension this chapter addresses is the
physical organisation of data within files. The choice between
row-oriented and columnar (column-oriented) storage can change the
I/O cost of an analytical query by one to two orders of magnitude,
making it one of the highest-leverage decisions in a data engineering
system.

%--------------------------------------------------------------------
\subsection{Row-Oriented Storage}
\label{subsec:row-oriented}
%--------------------------------------------------------------------

In a \term{row-oriented} (or \emph{N-ary storage model}, NSM) file
format, all fields of a single record are stored contiguously on disk:

\begin{center}
\setlength{\fboxsep}{4pt}
\fbox{\texttt{user\_id=1, name=``Alice'', age=24, score=91, country=``US'', ...}}\\[2pt]
\fbox{\texttt{user\_id=2, name=``Bob'',   age=31, score=78, country=``DE'', ...}}\\[2pt]
\fbox{\texttt{user\_id=3, name=``Carol'', age=27, score=88, country=``JP'', ...}}
\end{center}

Row-oriented storage is optimal for workloads that need to read or
write complete records: transactional databases (OLTP), log ingestion
pipelines, and streaming append operations. Writing a new record requires
a single sequential write. Reading a specific record by key requires
reading one contiguous block of bytes.

Row-oriented storage is \emph{suboptimal} for analytical workloads that
select a small number of columns from a wide table. A query
\texttt{SELECT score FROM users} must read every field of every row---
\texttt{user\_id}, \texttt{name}, \texttt{age}, \texttt{country},
and all other columns---even though only the \texttt{score} field is
needed. For a table with $c$ columns, a query selecting $k$ columns
reads $c/k$ times more data than necessary.

%--------------------------------------------------------------------
\subsection{Columnar Storage: The I/O Reduction Formula}
\label{subsec:columnar}
%--------------------------------------------------------------------

In a \term{columnar} (or \emph{decomposition storage model}, DSM) file
format, all values of a single column are stored contiguously on disk:

\begin{center}
\setlength{\fboxsep}{4pt}
\fbox{\texttt{user\_id: [1, 2, 3, ...]}} $\;\;\;$
\fbox{\texttt{name: [Alice, Bob, Carol, ...]}} $\;\;\;$
\fbox{\texttt{score: [91, 78, 88, ...]}}
\end{center}

For the same \texttt{SELECT score FROM users} query, a columnar store
reads \emph{only the score column chunk}---a contiguous block containing
\texttt{[91, 78, 88, ...]}, and nothing else. The \texttt{user\_id},
\texttt{name}, \texttt{age}, and other column chunks are not read at all.

\begin{defbox}{Columnar I/O Reduction}
Let a table have $c$ columns, $n$ rows, and $\bar{s}$ bytes per cell on
average. A query selecting $k$ columns:
\begin{align}
\text{Row store I/O}     &= n \cdot c \cdot \bar{s} \label{eq:row-io} \\
\text{Column store I/O}  &= n \cdot k \cdot \bar{s} \label{eq:col-io} \\
\text{I/O reduction}     &= \frac{k}{c} \label{eq:io-reduction}
\end{align}
For $k = 3$, $c = 100$: the columnar store reads only $3\%$ of the
data that a row store would read---a $33\times$ reduction.
\end{defbox}

Additional compression efficiency: because all values in a column chunk
have the same data type, they compress far better than a mixed row.
A column of \texttt{country} codes (``US'', ``DE'', ``JP'', ...) with
high repetition compresses 10--50$\times$ better than an entire row
containing a mix of integers, strings, and floats.

\begin{figure}[ht]
\centering
\begin{tikzpicture}[
    hdr/.style={rectangle, fill=DEBlue!80, text=white, draw=DEBlue,
                minimum width=1.5cm, minimum height=0.55cm,
                font=\tiny\bfseries\sffamily, align=center},
    rcell/.style={rectangle, draw=DEGray!40, fill=DEGray!8,
                  minimum width=1.5cm, minimum height=0.55cm,
                  font=\tiny, align=center},
    ccell/.style={rectangle, draw=DEGray!40, fill=DEGray!8,
                  minimum width=1.5cm, minimum height=0.55cm,
                  font=\tiny, align=center},
    hcell/.style={rectangle, draw=DEGreen!60, fill=DELightGreen,
                  minimum width=1.5cm, minimum height=0.55cm,
                  font=\tiny\bfseries, align=center},
    lbl/.style={font=\small\bfseries\sffamily, align=center}
  ]

  %--- Row Store (left) ---
  \node[lbl, text=DERed] at (-3.4, 4.1) {Row Store};
  \node[font=\tiny\sffamily\itshape, text=DEGray] at (-3.4, 3.7) {(CSV, JSON)};

  % Headers
  \foreach \col/\x in {user\_id/-5.3, name/-3.8, age/-2.3, score/-0.8}
    \node[hdr] at (\x, 3.2) {\col};

  % Row 1 (entire row highlighted red = all data read)
  \foreach \val/\x in {1/-5.3, Alice/-3.8, 24/-2.3, 91/-0.8}
    \node[rcell, fill=DERed!20, draw=DERed!50] at (\x, 2.65) {\val};
  % Row 2
  \foreach \val/\x in {2/-5.3, Bob/-3.8, 31/-2.3, 78/-0.8}
    \node[rcell, fill=DERed!20, draw=DERed!50] at (\x, 2.10) {\val};
  % Row 3
  \foreach \val/\x in {3/-5.3, Carol/-3.8, 27/-2.3, 88/-0.8}
    \node[rcell, fill=DERed!20, draw=DERed!50] at (\x, 1.55) {\val};
  % Row 4
  \foreach \val/\x in {4/-5.3, Dave/-3.8, 29/-2.3, 95/-0.8}
    \node[rcell, fill=DERed!20, draw=DERed!50] at (\x, 1.00) {\val};

  \draw[DERed!60, very thick, rounded corners=2pt]
       (-6.1, 0.7) rectangle (-0.1, 2.95);
  \node[font=\tiny\bfseries, text=DERed] at (-3.1, 0.4)
       {Query reads ALL 4 columns};

  %--- Column Store (right) ---
  \node[lbl, text=DEGreen!70!black] at (4.5, 4.1) {Column Store};
  \node[font=\tiny\sffamily\itshape, text=DEGray] at (4.5, 3.7) {(Parquet, ORC)};

  % Headers
  \foreach \col/\x in {user\_id/1.5, name/3.0, age/4.5, score/6.0}
    \node[hdr] at (\x, 3.2) {\col};

  % Column data stored contiguously
  \foreach \val/\y in {1/2.65, 2/2.10, 3/1.55, 4/1.00}
    \node[ccell] at (1.5, \y) {\val};
  \foreach \val/\y in {Alice/2.65, Bob/2.10, Carol/1.55, Dave/1.00}
    \node[ccell] at (3.0, \y) {\val};
  \foreach \val/\y in {24/2.65, 31/2.10, 27/1.55, 29/1.00}
    \node[ccell] at (4.5, \y) {\val};
  % score column highlighted (only one read)
  \foreach \val/\y in {91/2.65, 78/2.10, 88/1.55, 95/1.00}
    \node[hcell] at (6.0, \y) {\val};

  \draw[DEGreen!60, very thick, rounded corners=2pt]
       (5.2, 0.7) rectangle (6.8, 2.95);
  \node[font=\tiny\bfseries, text=DEGreen!70!black] at (4.5, 0.4)
       {Query reads ONLY score column};

  % Divider
  \draw[DEGray!30, dashed, thick] (0, 0.1) -- (0, 4.2);

\end{tikzpicture}
\caption{Row-oriented vs.\ columnar storage for a \texttt{SELECT score FROM users}
         query. The row store reads all four columns for every row (red box).
         The column store reads only the \texttt{score} column chunk (green box),
         skipping \texttt{user\_id}, \texttt{name}, and \texttt{age} entirely.
         For 100 columns with a 3-column query, the columnar store provides a
         $33\times$ I/O reduction.}
\label{fig:row-vs-column}
\end{figure}

%%====================================================================
\section{Big Data File Formats in Depth}
\label{sec:file-formats}
%%====================================================================

Three file formats dominate the modern Hadoop and cloud data lake
ecosystem: Apache Avro, Apache Parquet, and Apache ORC. Each represents
a different design point in the space of storage orientation, schema
management, compression efficiency, and processing engine compatibility.

%--------------------------------------------------------------------
\subsection{Apache Avro: Schema-Evolving Row Storage}
\label{subsec:avro}
%--------------------------------------------------------------------

Apache Avro is a row-oriented serialisation framework developed at
Yahoo! in 2009 and now an Apache top-level project. Avro stores each
record as a single contiguous binary blob, with the schema described
as a JSON document embedded in the file header. It is the preferred
format for data ingestion pipelines---particularly Kafka-to-HDFS
streams---for two reasons.

\textbf{Schema evolution.} As event schemas change over time (new
fields added, field names changed, types widened), Avro allows
\emph{reader schemas} and \emph{writer schemas} to differ according
to well-defined compatibility rules. A consumer reading data written
with an older schema can specify a new reader schema; Avro's binary
codec automatically:
\begin{itemize}
  \item Fills missing fields with their declared default values.
  \item Ignores fields present in the data but absent from the reader schema.
  \item Promotes numeric types (e.g., \texttt{int} to \texttt{long}).
\end{itemize}
This forward and backward compatibility makes Avro the standard for
streaming ingestion (Kafka) to long-term storage (HDFS), where the
schema evolves continuously as the application changes.

\textbf{Splittability.} An Avro file is divided into \term{blocks}
(not to be confused with HDFS blocks) separated by sync markers. A
Hadoop InputFormat can seek to a sync marker and start reading from
any block without decompressing from the beginning, making Avro files
natively splittable.

%--------------------------------------------------------------------
\subsection{Apache Parquet: Columnar Analytics Storage}
\label{subsec:parquet}
%--------------------------------------------------------------------

Apache Parquet, developed jointly by Twitter and Cloudera in 2013, is
the de facto standard for analytical data storage in the Hadoop
ecosystem and on cloud data lakes. Its internal organisation is
designed to minimise I/O for column-selective analytical queries.

\textbf{File structure.} A Parquet file is divided into \term{row
groups} (default $\sim$128\MB, aligning with HDFS block size). Each
row group is subdivided into \term{column chunks}---one per column.
Each column chunk is further divided into \term{pages} of approximately
1\MB\ each. Data within a page is optionally dictionary-encoded
(replacing repeated values with short integer codes) and then
compressed with a configurable codec (Snappy by default).

\begin{figure}[ht]
\centering
\begin{tikzpicture}[
    file/.style={rectangle, draw=DEBlue!50, fill=DELightBlue!20,
                 rounded corners=4pt, minimum width=8.5cm},
    rg/.style={rectangle, draw=DEBlue!40, fill=DEBlue!8,
               minimum width=8.0cm, minimum height=0.55cm,
               font=\small\bfseries\sffamily, align=center},
    cc/.style={rectangle, draw=DEGreen!50, fill=DELightGreen,
               minimum width=2.3cm, minimum height=0.5cm,
               font=\scriptsize\sffamily, align=center},
    pg/.style={rectangle, draw=DEAccent!60, fill=DELightAccent,
               minimum width=0.65cm, minimum height=0.4cm,
               font=\tiny, align=center},
    hdr/.style={rectangle, draw=DEBlue!70, fill=DEBlue!80, text=white,
                minimum width=8.0cm, minimum height=0.5cm,
                font=\scriptsize\bfseries\sffamily, align=center},
    ft/.style={rectangle, draw=DEAccent!80, fill=DEAccent!70, text=white,
               minimum width=8.0cm, minimum height=0.6cm,
               font=\scriptsize\bfseries\sffamily, align=center},
    lbl/.style={font=\scriptsize\sffamily, align=center}
  ]

  % File header
  \node[hdr] at (0, 7.2) {Magic bytes: PAR1};

  % Row Group 1
  \node[rg] at (0, 6.6) {Row Group 1 ($\approx$128 MB)};
  \node[cc] at (-2.6, 6.0) {col:user\_id};
  \node[cc] at (0,    6.0) {col:name};
  \node[cc] at (2.6,  6.0) {col:score};
  % Pages inside col:score
  \foreach \x in {2.0, 2.65, 3.3} {
    \node[pg] at (\x, 5.4) {P};
  }
  \node[font=\tiny\sffamily\itshape, text=DEGray] at (0, 5.4)
       {(each column chunk split into $\approx$1 MB pages)};

  % Row Group 2
  \node[rg] at (0, 4.8) {Row Group 2 ($\approx$128 MB)};
  \node[cc] at (-2.6, 4.2) {col:user\_id};
  \node[cc] at (0,    4.2) {col:name};
  \node[cc] at (2.6,  4.2) {col:score};

  % Footer
  \node[ft] at (0, 3.5) {File Footer};
  \node[lbl, text=DEAccent!80] at (0, 2.95)
       {\tiny Row group offsets \textbullet\ Column statistics (min/max/null count) per column per row group};
  \node[lbl, text=DEAccent!80] at (0, 2.55)
       {\tiny Schema \textbullet\ Encoding metadata \textbullet\ Compression codec};
  \node[hdr] at (0, 2.1) {Magic bytes: PAR1};

  % Predicate pushdown annotation
  \draw[DERed!60, very thick, dashed, rounded corners=3pt]
       (-4.4, 2.0) rectangle (4.4, 2.75);
  \node[font=\tiny\bfseries, text=DERed] at (5.5, 2.4)
       {Read footer FIRST};
  \node[font=\tiny\sffamily, text=DERed, align=left] at (5.5, 1.9)
       {skip row groups where\\col stats eliminate them};

\end{tikzpicture}
\caption{Parquet file internal structure. The footer contains per-column,
         per-row-group statistics (min, max, null count) that enable predicate
         pushdown: a query with \texttt{WHERE score > 90} reads the footer
         first, then skips all row groups where \texttt{max(score) $\leq$ 90},
         reading zero bytes from those row groups.}
\label{fig:parquet-layout}
\end{figure}

\textbf{Predicate pushdown.} The Parquet \term{file footer} stores
column-level statistics---minimum value, maximum value, and null
count---for every column chunk in every row group. A query engine
(Spark SQL, Hive on Tez, Presto) reads the footer first. If the
footer shows that the \texttt{score} column in Row Group 1 has
$\text{min}=60, \text{max}=85$, and the query filters
\texttt{WHERE score > 90}, the engine can skip the entire Row Group 1
without reading a single byte of its data. This \emph{predicate
pushdown} to the storage layer can reduce the physical I/O of a
selective query from terabytes to gigabytes.

\textbf{Dictionary encoding.} Before compression, Parquet applies
\term{dictionary encoding} to columns with low cardinality. A column
containing country codes with 200 distinct values out of 100 million
rows stores the 200 distinct values in a compact dictionary, then
stores each row's value as a small integer index into the dictionary.
This can reduce the size of a low-cardinality column by 80--95\%
before compression is even applied.

%--------------------------------------------------------------------
\subsection{Apache ORC: Hive-Optimised Columnar Storage}
\label{subsec:orc}
%--------------------------------------------------------------------

Apache ORC (Optimised Row Columnar) was developed at Hortonworks in
2013 specifically for Apache Hive and is the recommended format for
the Hive ecosystem. ORC and Parquet are architecturally similar---both
are columnar, self-describing, and support predicate pushdown---but
differ in several engineering choices.

ORC uses \term{stripes} (analogous to Parquet's row groups, default
250\MB) and stores a \term{file-level index} in the footer plus a
\term{stripe-level index} within each stripe. The stripe-level index
provides statistics at 10{,}000-row granularity within the stripe,
enabling finer-grained skipping than Parquet's row-group granularity.

ORC's default compression codec is Zlib (gzip), which achieves higher
compression ratios than Parquet's default Snappy at the cost of slower
decompression speed. For Hive workloads where query latency is measured
in minutes rather than seconds, the storage savings from Zlib typically
outweigh the decompression overhead. ORC also supports Hive's ACID
(atomicity, consistency, isolation, durability) transactions---a feature
Parquet does not natively support---enabling \texttt{UPDATE} and
\texttt{DELETE} operations on Hive tables.

\begin{table}[ht]
\centering
\caption{Comparison of the five major Hadoop file formats}
\label{tab:format-comparison}
\renewcommand{\arraystretch}{1.5}
\begin{tabular}{L{1.6cm} L{1.6cm} L{2.0cm} L{1.4cm} L{1.8cm} L{2.6cm}}
\toprule
\textbf{Format} & \textbf{Orient.} & \textbf{Schema} & \textbf{Split.} &
\textbf{Compress.} & \textbf{Primary use case} \\
\midrule
Text/CSV  & Row      & External        & With codec  & None (default) & Small data; portability \\
JSON      & Row      & Self (verbose)  & No (raw)    & External       & APIs; log ingestion \\
Avro      & Row      & Self (JSON)     & Yes         & Snappy/Deflate & Kafka ingest; schema evolution \\
Parquet   & Columnar & Self            & Yes         & Snappy/Gzip    & Spark, Presto, ML; analytics \\
ORC       & Columnar & Self            & Yes         & Zlib/Snappy    & Hive warehouse; ACID; archival \\
\bottomrule
\end{tabular}
\end{table}

%%====================================================================
\section{Compression Codecs and Splittability}
\label{sec:compression}
%%====================================================================

Compression reduces the physical size of data on disk and in transit,
lowering both storage costs and I/O time. In the Hadoop ecosystem,
choosing the right compression codec requires balancing three
dimensions: compression ratio, decompression speed, and---most
critically for distributed processing---\emph{splittability}.

%--------------------------------------------------------------------
\subsection{The Splittability Property}
\label{subsec:splittability}
%--------------------------------------------------------------------

\begin{defbox}{Splittability}
A file format is \textbf{splittable} if a Hadoop InputFormat can start
reading from any point in the file without decompressing from the
beginning. Splittability determines whether HDFS can assign multiple
independent Map tasks to different portions of the same file, enabling
full parallelism.
\end{defbox}

The practical consequence is dramatic. A 10\GB\ unsplittable Gzip CSV
file on HDFS is stored in 80 consecutive 128\MB\ blocks. HDFS distributes
these blocks across the cluster normally, but the MapReduce or Spark job
can assign only a \emph{single} task to this file: the task must read
all 80 blocks sequentially from start to finish, since decompression
must begin at byte 0 and proceed in order. No parallelism is possible.

A 10\GB\ Parquet+Snappy file with 128\MB\ row groups is similarly
stored in 80 blocks. But the Parquet InputFormat can start reading at
any row group boundary independently. The job assigns 80 parallel tasks,
each reading one row group from a local DataNode---achieving 80$\times$
parallelism and completing the job in $1/80$th the wall-clock time of
the single-task case.

\begin{figure}[ht]
\centering
\begin{tikzpicture}[
    blk/.style={rectangle, draw=DEGreen!60, fill=DELightGreen,
                minimum width=1.3cm, minimum height=0.65cm,
                font=\scriptsize\sffamily, align=center},
    gzblk/.style={rectangle, draw=DERed!60, fill=DERed!15,
                  minimum width=6.0cm, minimum height=0.65cm,
                  font=\scriptsize\sffamily, align=center},
    arr/.style={-{Stealth[length=4pt,width=4pt]}, thick},
    lbl/.style={font=\scriptsize\sffamily, align=center}
  ]

  % Parquet (splittable)
  \node[font=\small\bfseries\sffamily, text=DEGreen!70!black] at (3.0, 4.0)
       {Parquet + Snappy (splittable)};
  \node[blk] (rg1) at (0.8, 3.2) {Row Grp 1};
  \node[blk] (rg2) at (2.4, 3.2) {Row Grp 2};
  \node[blk] (rg3) at (4.0, 3.2) {Row Grp 3};
  \node[blk] (rg4) at (5.6, 3.2) {Row Grp 4};

  \draw[arr, DEGreen!70!black] (0.8, 2.7) -- (rg1.south) node[below, lbl]{\tiny Task 1};
  \draw[arr, DEGreen!70!black] (2.4, 2.7) -- (rg2.south) node[below, lbl]{\tiny Task 2};
  \draw[arr, DEGreen!70!black] (4.0, 2.7) -- (rg3.south) node[below, lbl]{\tiny Task 3};
  \draw[arr, DEGreen!70!black] (5.6, 2.7) -- (rg4.south) node[below, lbl]{\tiny Task 4};

  \node[font=\tiny\sffamily\itshape, text=DEGreen!70!black] at (3.2, 1.9)
       {4 tasks in parallel — full data locality possible};

  % Gzip CSV (not splittable)
  \node[font=\small\bfseries\sffamily, text=DERed] at (3.0, 1.3)
       {Gzip CSV (NOT splittable)};
  \node[gzblk] at (3.0, 0.7) {Single undivided compressed stream};

  \draw[arr, DERed!60] (3.0, 0.1) -- (3.0, 0.38)
       node[below, lbl, text=DERed] {\tiny Task 1 only --- no parallelism};

\end{tikzpicture}
\caption{Splittability comparison. Parquet's row-group boundaries allow
         $N$ tasks to process $N$ row groups in parallel. A Gzip-compressed
         CSV file has no internal boundaries the InputFormat can seek to,
         so the entire file is processed by a single task.}
\label{fig:splittability}
\end{figure}

%--------------------------------------------------------------------
\subsection{Compression Codec Comparison}
\label{subsec:codecs}
%--------------------------------------------------------------------

\begin{table}[ht]
\centering
\caption{Hadoop compression codecs: performance and splittability}
\label{tab:codecs}
\renewcommand{\arraystretch}{1.5}
\begin{tabular}{L{1.6cm} L{1.4cm} L{1.4cm} L{1.4cm} L{4.0cm}}
\toprule
\textbf{Codec} & \textbf{Speed} & \textbf{Ratio} & \textbf{Splittable?} &
\textbf{Recommended use} \\
\midrule
Snappy & Fastest & Medium   & No (but Parquet/Avro row groups are) & Parquet, Avro in Spark/Hive (default) \\
Gzip   & Slow    & High     & \textbf{No} & Avoid for large HDFS files; use Parquet+Snappy instead \\
LZO    & Fast    & Medium   & Yes (with index) & Legacy Hadoop; requires external index file \\
Bzip2  & Slowest & Highest  & Yes (natively) & Large plain-text files where I/O cost dominates \\
Zstd   & Fast    & High     & No (block-level) & ORC archival; Parquet cold storage \\
\bottomrule
\end{tabular}
\end{table}

\begin{caution}{The Gzip Anti-Pattern}
The most common performance anti-pattern in Hadoop data pipelines is
storing large datasets as Gzip-compressed CSV or JSON files.
Gzip's block format does not contain synchronisation markers that allow
random seeks; decompression must begin at byte 0. A 1~TB Gzip CSV file
processed by a 1{,}000-node Spark cluster is forced to run on
\emph{a single task}, using one core on one node, while the other
999 nodes sit idle. The correct alternative is \textbf{Parquet + Snappy}:
the Parquet format provides splittable boundaries (row groups) even
though Snappy itself is not splittable, and Snappy's very fast
decompression minimises CPU overhead.
\end{caution}

%%====================================================================
\section{I/O Optimisation Techniques}
\label{sec:io-opt}
%%====================================================================

Beyond format and codec selection, three additional query-time
optimisation techniques reduce the I/O cost of analytical workloads
on columnar storage systems.

%--------------------------------------------------------------------
\subsection{Predicate Pushdown}
\label{subsec:pred-pushdown}
%--------------------------------------------------------------------

\term{Predicate pushdown} moves filter conditions from the query engine
into the storage layer. In Parquet, the file footer stores column-level
statistics (min, max, null count) per row group. When a query contains
a filter predicate (e.g., \texttt{WHERE revenue > 1000000}), the query
engine evaluates the predicate against the footer statistics for each
row group \emph{before reading the row group's data}. Any row group for
which the predicate is definitely false (i.e., \texttt{max(revenue)}
$\leq 1{,}000{,}000$) is skipped entirely. Only row groups where the
predicate \emph{might} be true are read.

In practice, for a time-series table partitioned by date and sorted by
user ID within each partition, predicate pushdown on a user-ID filter
can skip 99\% of the data in large Parquet files, reducing a 1-hour
query to under 2 minutes.

%--------------------------------------------------------------------
\subsection{Column Pruning}
\label{subsec:col-pruning}
%--------------------------------------------------------------------

\term{Column pruning} (also called \emph{projection pushdown}) ensures
that only the columns required by the query are read from disk. In a
row-oriented format, this is impossible: all columns of a row are stored
together, so reading one requires reading all. In a columnar format, each
column is stored independently; the query engine passes the list of
required columns to the file reader, which opens and reads only those
column chunks. A query selecting 5 of 200 columns from a Parquet file
reads exactly $5/200 = 2.5\%$ of the storage bytes.

Apache Spark applies column pruning automatically for Parquet files
through the Catalyst query optimiser: the optimiser analyses the query's
column references before execution and passes only the referenced
columns to the Parquet reader. No manual configuration is required.

%--------------------------------------------------------------------
\subsection{Partition Pruning}
\label{subsec:partition-pruning-ch5}
%--------------------------------------------------------------------

\term{Partition pruning}, introduced in the context of Hive in
Chapter~\ref{ch:mapreduce}, operates at the file system (HDFS directory)
level rather than within a file. A table partitioned by \texttt{dt}
(date) stores each day's data in a separate HDFS subdirectory
(\texttt{/data/table/dt=2024-01-15/}). A query filtering on a specific
date reads only that subdirectory's files, skipping all other partitions
entirely. Partition pruning, predicate pushdown, and column pruning are
complementary and can be applied simultaneously:

\begin{enumerate}
  \item \textbf{Partition pruning} eliminates entire HDFS directories
        based on the partition key filter.
  \item \textbf{Predicate pushdown} eliminates row groups within the
        remaining Parquet files based on column statistics in the footer.
  \item \textbf{Column pruning} eliminates irrelevant column chunks
        within the remaining row groups.
\end{enumerate}

The combined effect can reduce the I/O of a targeted analytical query
from terabytes to megabytes---multiple orders of magnitude---purely
through optimisations at the storage layer, without any changes to the
query itself.

%%====================================================================
\section{Case Study: LinkedIn's Multi-Framework YARN Cluster (2013--2016)}
\label{sec:linkedin-case}
%%====================================================================

\begin{casestudy}{LinkedIn --- Consolidating Five Frameworks on One 5{,}000-Node YARN Cluster}

\textbf{Background.}
By 2013, LinkedIn was processing approximately 500~TB of data per day
for its professional networking platform---analysing member profiles,
connections, job postings, messages, and advertising signals to power
features used by over 300 million members. The company ran five distinct
data processing frameworks: MapReduce (for Hive ETL), Apache Spark
(for the People You May Know recommendation algorithm, PYMK), Apache
Samza (for real-time Kafka stream processing), Apache Pig (for ad
targeting feature engineering), and Apache Tez (for interactive
Hive-on-Tez queries). Before YARN, these frameworks ran on \emph{separate
siloed clusters}---three distinct Hadoop clusters, each configured and
operated independently.

\smallskip
\textbf{The silo problem.}
The siloed cluster architecture created two systemic inefficiencies.
First, each cluster was under-utilised: the MapReduce cluster averaged
45\% utilisation (heavily used during nightly batch runs, idle during
the day), while the Spark cluster was heavily utilised for the PYMK
job's 4-hour compute window and nearly idle the rest of the time.
Combining both clusters would have provided the same total capacity
at significantly higher average utilisation. Second, cross-cluster
data movement added 30--90 minutes of latency to pipelines that
needed data from multiple clusters: an ETL job on the MapReduce cluster
had to export results to HDFS and wait for the Spark cluster to ingest
them.

\smallskip
\textbf{The unified YARN architecture.}
In 2013, LinkedIn migrated to a unified 5{,}000-node YARN cluster using
the Capacity Scheduler with six queues:

\begin{center}
\begin{tabular}{lcc}
\toprule
\textbf{Queue}      & \textbf{Capacity} & \textbf{Primary workload} \\
\midrule
\texttt{production} & 40\%              & Hive ETL; Pig ad features \\
\texttt{spark}      & 25\%              & PYMK; Spark ML models \\
\texttt{samza}      & 15\%              & Kafka stream processing \\
\texttt{adhoc}      & 10\%              & Analyst Hive queries \\
\texttt{tez}        & 7\%               & Interactive Hive-on-Tez \\
\texttt{reserve}    & 3\%               & Overflow / burst capacity \\
\bottomrule
\end{tabular}
\end{center}

\smallskip
\textbf{The PYMK pipeline.}
LinkedIn's People You May Know (``Friend recommendations'') feature
illustrates how multiple frameworks share the YARN cluster in a single
end-to-end pipeline:
\begin{enumerate}
  \item \textbf{Ingest} (Samza, \texttt{samza} queue): a Kafka consumer
        reads member-profile-view events at approximately 2 million
        events per second, materialising them to HDFS in Avro format.
  \item \textbf{Feature store} (Hive, \texttt{production} queue): a
        nightly Hive ETL job joins first-, second-, and third-degree
        connection data for all 300 million members into a wide feature
        table stored in ORC+Zlib format (partitioned by member-ID bucket).
  \item \textbf{Graph model} (Spark GraphX, \texttt{spark} queue): a
        Spark job runs label propagation over the social graph to score
        approximately 10 billion candidate recommendation pairs per day.
        The PYMK Spark job used 600 executors each with 8~GB RAM and
        4 vCPUs---4{,}800~GB of RAM and 2{,}400 vCPUs drawn from the
        \texttt{spark} queue's 25\% allocation of the 5{,}000-node cluster.
  \item \textbf{Serve}: top-$k$ recommendations per member are written
        to Voldemort (LinkedIn's key-value store) for real-time serving
        at sub-10-millisecond latency.
\end{enumerate}

\smallskip
\textbf{Quantified outcomes.}
The migration from three siloed clusters to a unified YARN cluster
produced measurable improvements across every operational dimension:
\begin{itemize}
  \item \textbf{Utilisation}: average cluster utilisation rose from
        $\sim$45\% (across three separate clusters) to $\sim$78\%
        (on the unified cluster)---a 73\% increase in effective
        compute capacity from the same hardware.
  \item \textbf{PYMK latency}: the PYMK Spark job recompute time
        dropped from 18 hours (on a separate, undersized Hadoop 1.x
        MapReduce cluster) to 4 hours (on YARN with Spark GraphX),
        enabling daily rather than every-other-day recommendation
        refreshes.
  \item \textbf{Operational cost}: consolidation eliminated one full
        team's worth of cluster operations overhead and reduced
        inter-cluster data synchronisation latency by 90\%.
  \item \textbf{Peak containers}: at peak, the unified cluster ran
        over 400{,}000 simultaneous containers across all six queues.
\end{itemize}

\smallskip
\textbf{The architectural lesson.}
LinkedIn's deployment established a principle that has since become an
industry standard: \emph{the cluster is the computer, and the resource
manager is its operating system}. Just as an operating system allows
multiple processes with different CPU and memory profiles to share
hardware fairly, YARN allows multiple data processing frameworks with
different resource profiles to share a cluster fairly. The Capacity
Scheduler's queue hierarchy provides the ``process priority'' mechanism;
the Container is the ``process memory allocation''; the NodeManager's
cgroup-based isolation is the ``memory protection''. This analogy---
cluster OS / cluster kernel---was explicitly used by the YARN paper's
authors and has guided the design of Kubernetes, Mesos, and every
subsequent cluster resource manager.
\end{casestudy}

%%====================================================================
\section*{Chapter Summary}
\label{sec:ch5-summary}
%%====================================================================

\begin{itemize}
  \item Hadoop 1.x's \textbf{JobTracker} bottleneck had three causes:
        framework lock-in (MapReduce only), rigid fixed-size slots,
        and single-process scalability limits ($\approx$4{,}000 nodes).

  \item \textbf{YARN} separates resource management into three
        components: \textbf{ResourceManager} (global scheduler),
        \textbf{NodeManager} (per-node agent), and
        \textbf{ApplicationMaster} (per-job framework logic).
        A \textbf{Container} (\texttt{\{node, vCores, memory\}})
        replaces the fixed slot.

  \item YARN's three schedulers: \textbf{FIFO} (single queue),
        \textbf{Capacity} (guaranteed multi-queue allocation),
        and \textbf{Fair} (equal sharing over time). \textbf{DRF}
        extends fairness to multiple resource dimensions.

  \item \textbf{Row-oriented} formats (CSV, JSON, Avro) store all
        fields of a record contiguously. \textbf{Columnar} formats
        (Parquet, ORC) store all values of a column contiguously.
        For a $k$-column query on a $c$-column table, columnar I/O
        is reduced by factor $k/c$.

  \item \textbf{Avro}: row-oriented, schema evolution (reader/writer
        schemas), splittable, Kafka-native. \textbf{Parquet}:
        columnar, row groups, footer statistics for predicate pushdown,
        Spark/Presto native. \textbf{ORC}: columnar, stripe-level
        index, Zlib compression, Hive ACID support.

  \item A file is \textbf{splittable} if multiple tasks can read
        different ranges independently. Gzip CSV is not splittable
        (one task only). Parquet/Avro blocks are splittable even if
        the codec (Snappy) is not.

  \item The three I/O optimisation layers apply in order:
        \textbf{partition pruning} (skip directories), \textbf{predicate
        pushdown} (skip row groups via footer statistics),
        \textbf{column pruning} (read only required column chunks).
\end{itemize}

%%====================================================================
\section*{Exercises}
\label{sec:ch5-exercises}
%%====================================================================

\exercisesection{Short Questions}

\sq What is the YARN ResourceManager's Scheduler component? What
information does the Scheduler use to make allocation decisions, and
what information does it deliberately \emph{not} know?

\sq Explain the role of the ApplicationMaster (AM) in YARN. Why does
one AM per job---rather than a centralised AM for all jobs---improve
the scalability of the ResourceManager?

\sq What is a YARN Container? Write the resource specification for a
container that would host a Spark executor requiring 8 CPU cores and
24~GB of RAM on node \texttt{dn-07.cluster.internal}.

\sq Explain the difference between YARN's Capacity Scheduler and Fair
Scheduler. Give a production scenario where the Capacity Scheduler is
the better choice, and a scenario where the Fair Scheduler is better.

\sq What is Dominant Resource Fairness (DRF)? Give an example where
allocating equal CPU across two jobs would result in an unfair outcome,
and explain how DRF resolves it.

\sq A YARN NodeManager stops sending heartbeats to the ResourceManager.
Trace the sequence of events that occurs: which component detects the
failure, what happens to the containers on that NodeManager, and which
component decides what to do next?

\sq What is the difference between the \emph{write schema} and the
\emph{read schema} in Apache Avro? Give a concrete example where they
differ and explain how Avro handles the mismatch.

\sq Explain predicate pushdown in Apache Parquet. What information is
stored in the Parquet file footer, and how does a query engine use it
to avoid reading data?

\sq What is the splittability property for a Hadoop file format? A
10\GB\ file on a 100-node cluster is stored as Gzip-compressed CSV.
How many Map tasks does MapReduce assign to this file, and why?

\sq Compare Parquet and ORC on three dimensions: primary processing
engine, default compression codec, and ACID transaction support.

\sq A data engineer proposes storing a 500\TB\ analytics warehouse as
raw JSON files (one JSON object per line) in HDFS. Identify three
specific performance problems with this proposal and propose a better
alternative.

\sq Explain column pruning (projection pushdown). For a 200-column
Parquet table, a query references 6 columns. What fraction of the
physical storage does the Parquet reader access?

\sq What does it mean for a columnar format to use \emph{dictionary
encoding}? For a column containing the values ``US'', ``GB'', ``DE'',
``FR'' repeated across 500 million rows, explain how dictionary encoding
reduces the column chunk size.

\sq YARN's ResourceManager HA uses ZooKeeper for failover. What role
does ZooKeeper play, and what happens to running applications during
the 30--60 second RM failover window?

\bigskip
\exercisesection{Analytical Problems}

\ap \textbf{YARN Capacity Planning.}
An e-commerce company operates a 2{,}000-node YARN cluster. Each node
has 64~GB RAM and 16 vCPUs. The Capacity Scheduler is configured with
three queues: \texttt{batch} (50\%), \texttt{streaming} (30\%),
\texttt{adhoc} (20\%). Container sizes: batch MapReduce mapper:
\{4~GB, 2~vCPU\}; streaming Samza task: \{2~GB, 1~vCPU\}; adhoc Spark
executor: \{8~GB, 4~vCPU\}.

\begin{enumerate}[label=(\alph*)]
  \item Calculate the total cluster RAM and vCPU capacity.
  \item For each queue, calculate the maximum number of concurrent
        containers, and identify whether the bottleneck resource is
        RAM or vCPU.
  \item A batch job needs 5{,}000 mapper containers simultaneously.
        How many batch nodes can serve this simultaneously? If demand
        exceeds the \texttt{batch} queue's allocation, describe what
        happens under (i) Capacity Scheduler with preemption and
        (ii) Capacity Scheduler without preemption.
  \item The streaming team submits a Samza job with 3{,}000 tasks,
        each requiring \{2~GB, 1~vCPU\}. Calculate the dominant
        resource for the streaming queue using DRF.
\end{enumerate}

\ap \textbf{File Format I/O Comparison.}
A retail company stores 5 years of transaction records in HDFS. The
table has 80 columns, 10 billion rows, and an average cell size of
12 bytes. A weekly analytics query selects 4 columns, filters on a
date range covering 7 days out of 365, and has a predicate on one
column that eliminates 90\% of row groups.

\begin{enumerate}[label=(\alph*)]
  \item Calculate the total uncompressed table size in terabytes.
  \item If stored as CSV, estimate the I/O for the weekly query
        (no partition pruning, no predicate pushdown).
  \item If stored as Parquet + Snappy, partitioned by date, with
        predicate pushdown: apply partition pruning (7/365 days),
        then predicate pushdown (90\% of remaining row groups skipped),
        then column pruning (4/80 columns). Calculate the remaining I/O.
  \item Express the I/O reduction as a ratio. If each I/O unit
        (TB read) costs the cluster 5 minutes of compute time on
        a 100-node cluster, estimate the query time saving.
\end{enumerate}

\ap \textbf{Splittability and Parallelism Analysis.}
A genomics pipeline ingests 2~TB of DNA sequencing data per day.
Three format options are proposed:

\begin{itemize}
  \item \textbf{Option A}: single 2~TB Gzip-compressed FASTQ file
  \item \textbf{Option B}: 16{,}000 Gzip-compressed FASTQ files of
        128~MB each
  \item \textbf{Option C}: Parquet files with 128~MB row groups
        (2~TB total across $\approx$16{,}000 row groups)
\end{itemize}

\begin{enumerate}[label=(\alph*)]
  \item For each option, calculate the number of Map tasks HDFS assigns,
        and explain whether data-local execution is achievable.
  \item Option B produces 16{,}000 files. What is the NameNode memory
        overhead (assume 300 bytes per file + block pair, and each
        file fits in one 128~MB block)?
  \item Option C uses a single logical Parquet file of 2~TB divided
        into 16{,}000 row groups. Assuming perfect data locality, how many
        Map tasks run in parallel if the cluster has 500 nodes with
        40 task slots each?
  \item Rank the three options on: (i) MapReduce parallelism, (ii)
        NameNode memory pressure, (iii) engineering simplicity.
        Recommend one option for the production pipeline.
\end{enumerate}

\ap \textbf{YARN Failure Cascade Analysis.}
A 500-node YARN cluster is running a Spark job with 400 executor
containers distributed across 400 nodes. Each executor holds 50~GB of
cached RDD data (4 partitions of 12.5~GB each). A power distribution
unit trip simultaneously takes down 50 nodes.

\begin{enumerate}[label=(\alph*)]
  \item How many executor containers are immediately lost?
  \item Which YARN components detect the node failure, and in what order?
  \item The Spark ApplicationMaster requests 50 replacement executors.
        If the remaining 450 nodes each have 4 available executor slots,
        how quickly can the RM satisfy the request (assume 500 ms
        scheduling cycle)?
  \item The replacement executors must recompute the 50 lost executors'
        cached RDD partitions (200 partitions, 12.5~GB each, at
        100~MB/s recompute throughput). Estimate the total recompute time.
  \item What data engineering design change would eliminate the need
        for full RDD recomputation after executor failure?
\end{enumerate}

\bigskip
\exercisesection{Design Problems}

\dpr \textbf{Data Lake Format Strategy for a Financial Services Platform.}
A global investment bank ingests 15 sources of market data into a central
data lake: real-time tick data (1 billion trades/day, each trade:
\{symbol, price, volume, timestamp, exchange, side\}), reference data
(static security metadata: 2 million instruments, 120 columns),
and regulatory reporting data (daily settlement reports, 80 columns,
7-year retention requirement).

Design the format and compression strategy for each data tier. Your
design must specify:
\begin{enumerate}[label=(\alph*)]
  \item The HDFS file format and compression codec for each of the
        three source types, with explicit justification for each choice.
  \item The partitioning scheme for tick data (which columns to partition
        on, at what granularity, and why).
  \item The YARN queue configuration for three consumer teams: algorithmic
        trading analytics (latency-sensitive, needs $<$5-minute query
        response), risk analytics (batch, 2-hour window nightly), and
        compliance reporting (weekly, long-running, low priority).
  \item How you would handle schema evolution if the tick data schema
        adds three new fields for a new exchange that emits extended
        market data.
  \item The estimated storage reduction from converting raw CSV tick
        data (at 35 bytes per event on average) to Parquet+Snappy
        (assuming 6:1 compression ratio on this workload) for one
        year of data.
\end{enumerate}

\dpr \textbf{YARN Migration from Three Siloed Clusters.}
A media company runs three separate Hadoop clusters: Cluster A
(1{,}000 nodes, MapReduce for nightly content ETL), Cluster B
(500 nodes, Spark for recommendation models), Cluster C
(300 nodes, Kafka/Samza for real-time user event processing).
Average utilisation: Cluster A 40\%, Cluster B 35\%, Cluster C 55\%.
The company plans to consolidate into a single YARN cluster.

Design the consolidation. Your design must address:
\begin{enumerate}[label=(\alph*)]
  \item Total node count for the unified cluster, given a target
        average utilisation of 70\% (justify your sizing).
  \item Capacity Scheduler queue configuration: queue names, capacity
        percentages, and maximum capacities for each team's workload.
  \item The migration sequence: which workload migrates first, how
        parallel operation during the transition is managed, and how
        you validate that the migrated workload performs correctly.
  \item How preemption should be configured between the nightly ETL
        jobs (MapReduce, 8-hour window) and the recommendation Spark
        jobs (4-hour window, started 2 hours before ETL completes).
  \item The rollback plan if the unified cluster shows unexpected
        performance degradation for the Samza real-time workload.
\end{enumerate}

\bigskip
\exercisesection{Programming Exercises}

\pe \textbf{YARN Capacity Planning Calculator (Python).}
Implement a YARN capacity planning tool that models a multi-queue
cluster, computes maximum concurrent containers per queue, identifies
the bottleneck resource (RAM vs.\ vCPU), and detects over-subscription.

\begin{lstlisting}[style=pyspark, caption={PE5.1: YARN capacity planning calculator}]
from dataclasses import dataclass

@dataclass
class NodeSpec:
    count:   int
    ram_gb:  int
    vcpus:   int

@dataclass
class QueueSpec:
    name:         str
    capacity_pct: float    # fraction of cluster [0, 1]
    container_ram_gb: int
    container_vcpus:  int

def plan_cluster(node: NodeSpec, queues: list[QueueSpec]):
    total_ram_gb = node.count * node.ram_gb
    total_vcpus  = node.count * node.vcpus

    print(f"Cluster: {node.count} nodes x {node.ram_gb} GB / {node.vcpus} vCPU")
    print(f"Total:   {total_ram_gb:,} GB RAM | {total_vcpus:,} vCPUs")
    print()
    print(f"{'Queue':<14} {'Cap':>5} {'RAM alloc':>10} {'CPU alloc':>10} "
          f"{'Max conts':>10} {'Bottleneck':>12}")
    print("-" * 65)

    for q in queues:
        q_ram_gb = total_ram_gb  * q.capacity_pct
        q_vcpus  = total_vcpus   * q.capacity_pct
        by_ram   = int(q_ram_gb  / q.container_ram_gb)
        by_cpu   = int(q_vcpus   / q.container_vcpus)
        max_c    = min(by_ram, by_cpu)
        bottle   = "RAM" if by_ram < by_cpu else "vCPU"
        print(f"{q.name:<14} {q.capacity_pct*100:>4.0f}% "
              f"{q_ram_gb:>9,.0f} GB {q_vcpus:>9,.0f}    "
              f"{max_c:>10,} {bottle:>12}")

    # Detect container spec errors: container bigger than node
    print()
    for q in queues:
        if q.container_ram_gb > node.ram_gb:
            print(f"[ERROR] Queue {q.name}: container RAM {q.container_ram_gb} GB "
                  f"> node RAM {node.ram_gb} GB")
        if q.container_vcpus > node.vcpus:
            print(f"[ERROR] Queue {q.name}: container vCPUs {q.container_vcpus} "
                  f"> node vCPUs {node.vcpus}")

# LinkedIn 5,000-node cluster model
node = NodeSpec(count=5_000, ram_gb=64, vcpus=16)

queues = [
    QueueSpec("production", 0.40, container_ram_gb=4,  container_vcpus=2),
    QueueSpec("spark",      0.25, container_ram_gb=8,  container_vcpus=4),
    QueueSpec("samza",      0.15, container_ram_gb=2,  container_vcpus=1),
    QueueSpec("adhoc",      0.10, container_ram_gb=4,  container_vcpus=2),
    QueueSpec("tez",        0.07, container_ram_gb=4,  container_vcpus=2),
    QueueSpec("reserve",    0.03, container_ram_gb=4,  container_vcpus=2),
]

plan_cluster(node, queues)
\end{lstlisting}

\pe \textbf{Parquet Layout and Predicate Pushdown Simulator (Python).}
Simulate the Parquet file structure---row groups, column chunks,
footer statistics---and implement predicate pushdown to skip row groups
for a simple filter predicate.

\begin{lstlisting}[style=pyspark, caption={PE5.2: Parquet predicate pushdown simulator}]
import random
import math

random.seed(42)

ROW_GROUP_ROWS = 1_000        # rows per row group (small for demo)
TOTAL_ROWS     = 10_000       # total table rows
COLUMNS        = ["user_id", "name", "country", "age", "score", "revenue"]
BYTES_PER_CELL = 8            # average bytes per cell

def generate_row_group(rg_id, start_row, rows):
    """Simulate one Parquet row group with column stats."""
    data = {
        "user_id": list(range(start_row, start_row + rows)),
        "score":   [random.randint(40, 100) for _ in range(rows)],
        "revenue": [random.randint(0, 1_000_000) for _ in range(rows)],
        "country": [random.choice(["US","GB","DE","JP","FR"]) for _ in range(rows)],
        "age":     [random.randint(18, 80) for _ in range(rows)],
        "name":    [f"User_{i}" for i in range(start_row, start_row + rows)],
    }
    # Build footer stats per column
    stats = {}
    for col, vals in data.items():
        if isinstance(vals[0], int):
            stats[col] = {"min": min(vals), "max": max(vals),
                          "null_count": 0}
        else:
            stats[col] = {"min": min(vals), "max": max(vals),
                          "null_count": 0}
    return {"rg_id": rg_id, "row_count": rows, "data": data, "stats": stats}

# Build Parquet file (footer + row groups)
num_rg  = math.ceil(TOTAL_ROWS / ROW_GROUP_ROWS)
rg_list = [generate_row_group(i, i * ROW_GROUP_ROWS,
                               min(ROW_GROUP_ROWS,
                                   TOTAL_ROWS - i * ROW_GROUP_ROWS))
           for i in range(num_rg)]

footer = {"num_row_groups": num_rg,
          "schema": COLUMNS,
          "row_group_stats": [rg["stats"] for rg in rg_list]}

def parquet_scan(row_groups, footer, predicate_col, predicate_op, threshold,
                 select_cols):
    """
    Simulate Parquet scan with:
      - predicate pushdown (skip row groups via footer stats)
      - column pruning (read only select_cols)
    """
    total_rg = len(row_groups)
    rg_read  = 0
    rows_returned = 0

    for i, rg in enumerate(row_groups):
        stats = footer["row_group_stats"][i]
        col_stat = stats.get(predicate_col, {})

        # Predicate pushdown: can we skip this row group?
        if predicate_op == ">" and col_stat.get("max", float("inf")) <= threshold:
            continue    # definitely no rows satisfy predicate -> skip
        if predicate_op == "<" and col_stat.get("min", float("-inf")) >= threshold:
            continue

        rg_read += 1
        # Column pruning: only access select_cols
        for row_i in range(rg["row_count"]):
            row_val = rg["data"][predicate_col][row_i]
            if predicate_op == ">" and row_val > threshold:
                rows_returned += 1
            elif predicate_op == "<" and row_val < threshold:
                rows_returned += 1

    # I/O calculation
    full_scan_bytes  = total_rg * ROW_GROUP_ROWS * len(COLUMNS)  * BYTES_PER_CELL
    actual_col_bytes = rg_read  * ROW_GROUP_ROWS * len(select_cols) * BYTES_PER_CELL
    io_reduction_pct = (1 - actual_col_bytes / full_scan_bytes) * 100

    print(f"Query: SELECT {select_cols} WHERE {predicate_col} {predicate_op} {threshold}")
    print(f"  Row groups total : {total_rg}")
    print(f"  Row groups read  : {rg_read}  (skipped {total_rg - rg_read} via pushdown)")
    print(f"  Rows returned    : {rows_returned:,}")
    print(f"  Full scan bytes  : {full_scan_bytes:,}")
    print(f"  Actual I/O bytes : {actual_col_bytes:,}")
    print(f"  I/O reduction    : {io_reduction_pct:.1f}%")

# Query: SELECT user_id, score WHERE score > 90
parquet_scan(rg_list, footer,
             predicate_col="score", predicate_op=">", threshold=90,
             select_cols=["user_id", "score"])
print()
# Query: SELECT user_id, revenue WHERE revenue > 800000
parquet_scan(rg_list, footer,
             predicate_col="revenue", predicate_op=">", threshold=800_000,
             select_cols=["user_id", "revenue"])
\end{lstlisting}

\pe \textbf{DRF Fairness Simulation (Python).}
Simulate Dominant Resource Fairness for a cluster with two resource
dimensions (CPU and RAM) and three competing jobs, showing how DRF
equalises dominant shares.

\begin{lstlisting}[style=pyspark, caption={PE5.3: Dominant Resource Fairness (DRF) simulator}]
def drf_allocate(cluster_cpu, cluster_ram_gb, jobs, max_rounds=20):
    """
    Simulate DRF allocation.
    Each job specifies: name, cpu_per_task, ram_per_task_gb, max_tasks
    DRF grants one task at a time to the job with the lowest dominant share.
    """
    allocated = {j["name"]: {"tasks": 0, "cpu": 0.0, "ram": 0.0}
                 for j in jobs}
    used_cpu = 0.0
    used_ram = 0.0

    print(f"Cluster: {cluster_cpu} CPUs | {cluster_ram_gb} GB RAM")
    print(f"\n{'Round':>5}  {'Job':>12}  {'Dominant':>10}  "
          f"{'Dom share':>10}  {'Tasks':>7}  {'CPU':>7}  {'RAM':>7}")
    print("-" * 68)

    for rnd in range(1, max_rounds + 1):
        # Compute dominant share for each job
        dom_shares = {}
        for j in jobs:
            a = allocated[j["name"]]
            if a["tasks"] >= j["max_tasks"]:
                dom_shares[j["name"]] = float("inf")
                continue
            cpu_share = a["cpu"] / cluster_cpu if cluster_cpu > 0 else 0
            ram_share = a["ram"] / cluster_ram_gb if cluster_ram_gb > 0 else 0
            dom_res   = "cpu" if cpu_share >= ram_share else "ram"
            dom_shares[j["name"]] = (max(cpu_share, ram_share), dom_res)

        # Select job with minimum dominant share
        eligible = [j for j in jobs
                    if isinstance(dom_shares[j["name"]], tuple)
                    and used_cpu + j["cpu_per_task"] <= cluster_cpu
                    and used_ram + j["ram_per_task_gb"] <= cluster_ram_gb]
        if not eligible:
            break

        winner = min(eligible, key=lambda j: dom_shares[j["name"]][0])
        w_name = winner["name"]
        dom_val, dom_res = dom_shares[w_name]

        allocated[w_name]["tasks"] += 1
        allocated[w_name]["cpu"]   += winner["cpu_per_task"]
        allocated[w_name]["ram"]   += winner["ram_per_task_gb"]
        used_cpu += winner["cpu_per_task"]
        used_ram += winner["ram_per_task_gb"]

        print(f"{rnd:>5}  {w_name:>12}  {dom_res:>10}  "
              f"{dom_val*100:>9.1f}%  "
              f"{allocated[w_name]['tasks']:>7}  "
              f"{allocated[w_name]['cpu']:>7.1f}  "
              f"{allocated[w_name]['ram']:>7.1f}")

    print("\nFinal allocations:")
    for j in jobs:
        a = allocated[j["name"]]
        cpu_share = a["cpu"] / cluster_cpu * 100
        ram_share = a["ram"] / cluster_ram_gb * 100
        print(f"  {j['name']:<14}: {a['tasks']} tasks | "
              f"CPU: {a['cpu']:.1f} ({cpu_share:.1f}%) | "
              f"RAM: {a['ram']:.1f} GB ({ram_share:.1f}%)")

# 3 jobs with different resource profiles
jobs = [
    {"name": "MapReduce",  "cpu_per_task": 2,  "ram_per_task_gb": 4,   "max_tasks": 8},
    {"name": "Spark",      "cpu_per_task": 4,  "ram_per_task_gb": 16,  "max_tasks": 6},
    {"name": "Samza",      "cpu_per_task": 1,  "ram_per_task_gb": 2,   "max_tasks": 12},
]

drf_allocate(cluster_cpu=32, cluster_ram_gb=128, jobs=jobs, max_rounds=20)
\end{lstlisting}
